\documentclass[letterpaper,10pt]{article}

% ----------------------------
% Paquetes
% ----------------------------
\usepackage{geometry}
\geometry{
  left=1cm,
  right=1cm,
  top=1.5cm,
  bottom=0.8cm,
  headheight=12pt
}

\usepackage[utf8]{inputenc} \usepackage[T1]{fontenc} \usepackage{geometry, graphicx, amsmath, amssymb, hyperref, fancyhdr, parskip, caption, xcolor, wrapfig, float, titlesec, setspace}

\usepackage[nopatch=footnote]{microtype}

\usepackage{csquotes}

\usepackage[backend=biber,style=apa]{biblatex}
\addbibresource{references.bib}

% ----------------------------
% Formato del documento
% ----------------------------
\setlength{\parskip}{0.4em}
\setlength{\parindent}{0.8em}
\setstretch{1} % interlineado compacto

\hypersetup{
    colorlinks = true,
    linkcolor = blue,
    citecolor = red,
    urlcolor  = blue
}

% ----------------------------
% Encabezado y pie de página
% ----------------------------
\pagestyle{fancy}
\fancyhf{}
\lhead{\scriptsize Learning the theory of Monte Carlo methods}
\rhead{\scriptsize Page \thepage}
\renewcommand{\headrulewidth}{0.4pt}

% ----------------------------
% Formato de títulos
% ----------------------------
\titleformat{\section}
  {\normalsize\bfseries}{\thesection.}{0.4em}{}

\titleformat{\subsection}
  {\small\bfseries}{\thesubsection.}{0.4em}{}

\titlespacing*{\section}{0pt}{0.8em}{0.3em}
\titlespacing*{\subsection}{0pt}{0.6em}{0.2em}

% ----------------------------
% Portada
% ----------------------------
\begin{document}

\begin{titlepage}
    \centering
    \vspace*{2cm}
  
    % ---- Main title ----
    {\huge\bfseries Learning the theory of Monte Carlo methods}\\[0.3cm]
    
    % ---- Sub title ----
    {\Large Based on Simon Sirca - Probability for physicists}\\[1cm]

    {\normalsize \textbf{Estiven Castrillon Alzate, \href{mailto:estiven.castrillon2@udea.edu.co}{\texttt{estiven.castrillon2@udea.edu.co}}}}\\
    {\normalsize \textbf{Date:} \today}\\[2cm]
\end{titlepage}

\section{Introduction}

The Monte-Carlo method is a tool used widely in computer simulations in experimental physics. Here, I will study this mathematical and computational technique to enhance my understanding of simulations in High Energy Physics (HEP). The main reference for this learning process is the book \textit{Probability for Physicists}~\parencite{Sirca_2016}. 

The Monte-Carlo method considered in this work is the acceptance-rejection method, a technique used to generate random samples from a probability distribution that may be difficult to sample from directly. The method works by generating random samples from a proposal distribution and accepting or rejecting those samples according to a criterion that guarantees that the accepted samples follow the desired target distribution~\parencite{Nakade_2025}.

\subsection{Acceptance-Rejection Method}

Let $p(x)$ be a probability density function (PDF) defined on an interval $[a,b]$, from which we want to sample. Suppose that there exists a constant $M > 0$ such that
\begin{equation}
    p(x) \leq M, \quad \forall x \in [a,b].
\end{equation}

The acceptance-rejection method is based on sampling uniformly from the rectangular region defined by $[a,b] \times [0,M]$, and retaining only those points that lie under the curve $y = p(x)$. According to the fundamental theorem of the method, the marginal distribution of the accepted $x$ values follows the target distribution $p(x)$~\parencite{Nakade_2025}.

\subsection{Algorithm}

The basic acceptance-rejection algorithm can be summarized as follows:

\begin{enumerate}
    \item \textbf{Generate a proposal:} Sample 
    \begin{equation}
        x \sim \mathcal{U}(a,b).
    \end{equation}
    
    \item \textbf{Generate a test value:} Sample 
    \begin{equation}
        y \sim \mathcal{U}(0,M).
    \end{equation}
    
    \item \textbf{Acceptance criterion:} Accept $x$ if
    \begin{equation}
        y \leq p(x),
    \end{equation}
    otherwise reject the sample and return to step 1.
\end{enumerate}

This procedure ensures that the accepted samples are distributed according to $p(x)$.

\subsection{Extension with Proposal Distributions}

A more general version of the acceptance-rejection method introduces a proposal distribution $g(x)$ and a constant $M$ such that
\begin{equation}
    p(x) \leq M g(x), \quad \forall x.
\end{equation}

The algorithm is then modified as follows:

\begin{enumerate}
    \item Sample 
    \begin{equation}
        x \sim g(x).
    \end{equation}
    
    \item Sample 
    \begin{equation}
        u \sim \mathcal{U}(0,1).
    \end{equation}
    
    \item Accept $x$ if
    \begin{equation}
        u \leq \frac{p(x)}{M g(x)},
    \end{equation}
    otherwise reject and repeat.
\end{enumerate}

This formulation is more efficient when $g(x)$ approximates the shape of $p(x)$, reducing the number of rejected samples.

\subsection{Physical Interpretation}

From a physical perspective, the acceptance-rejection method allows us to simulate random variables distributed according to a given differential rate. In the context of particle physics, if the differential decay rate $\frac{d\Gamma}{dx}$ is known, it can be interpreted as a probability density (up to normalization), and the method enables the generation of synthetic decay events that reproduce the expected physical distribution.

\printbibliography
\end{document}